Grounded VideoQA requires models not only to localize relevant temporal evidence but also to reason over that evidence to predict the correct answer. We identify a key challenge in multiple-choice grounded VideoQA: the mismatch between answer-oriented candidate text and evidence-oriented retrieval. \emph{EviR} addresses this issue by decomposing grounded prediction into query reformulation, temporal retrieval, evidence validation, and answer generation, enabling more faithful evidence selection before final reasoning.

Experiments on NExT-GQA and ReXTime demonstrate that this design consistently improves grounding-aware performance, with the largest gains observed on stricter metrics that require precise evidence localization. Ablation results further reveal that the Refiner is essential for improving retrieval alignment, while the Generator remains critical for resolving ambiguity among answer candidates even when strong evidence is available. These findings highlight that high-quality evidence localization alone is not sufficient for grounded VideoQA, and that explicit reasoning over validated evidence is necessary for reliable prediction.

Overall, the results suggest that grounded VideoQA benefits from explicitly decoupling evidence acquisition from answer prediction, rather than relying on answer candidates as direct retrieval queries. A limitation of the current framework is its reliance on a single top-ranked evidence clip, which may be insufficient for more complex scenarios involving multiple relevant events or compositional reasoning. Future work includes extending \emph{EviR} to open-ended grounded VideoQA and exploring multi-evidence aggregation for more expressive and robust temporal reasoning.
